% ============================================================
% CHAPTER 3: Planning and Architecture
% (Following Tarek's Chapter 2 structure — Requirements, UML, Architecture, Tools)
% ============================================================

\chapter{Planning and Architecture}

\section{Introduction}

This chapter presents the sprint zero, which represents the first step toward the realization of our project. First, we identify the actors of our application. Then, we list the functional and non-functional requirements of our system and detail the work through the methodology chosen in the previous chapter. Finally, we provide a brief overview of the hardware, technologies, and programming languages used to set up the working environment.

% ============================================================
\section{Requirements Capture}
% ============================================================

\subsection{Identification of Actors}

The application involves five types of actors, each with specific responsibilities according to their role in the system.

\vspace{0.3cm}
\noindent\textbf{Developer} \\
The developer is the primary end-user who writes infrastructure code (Terraform, Dockerfile, Kubernetes YAML). Developers interact indirectly with the service through the CI/CD pipeline: they push code to Git, which triggers automated scanning. When violations are detected, developers receive detailed feedback including line numbers, violation messages, and remediation recommendations.

\vspace{0.3cm}
\noindent\textbf{DevOps Engineer} \\
The DevOps engineer manages the CI/CD pipeline and deployment infrastructure. This actor configures Jenkins pipelines, monitors the service metrics dashboard, reviews scan history, and adjusts the block threshold settings. The DevOps engineer interacts directly with the service through the metrics and history endpoints.

\vspace{0.3cm}
\noindent\textbf{Security Engineer} \\
The security engineer defines and maintains the security policies enforced by the service. This actor configures policies via \texttt{policies\_config.yaml}, adjusts severity levels, enables or disables specific rules, and monitors violation trends to identify areas requiring developer education.

\vspace{0.3cm}
\noindent\textbf{Jenkins / CI System (Secondary Actor)} \\
Jenkins is the automation orchestrator that detects code changes, triggers security scans by calling the \texttt{POST /analyze} endpoint, parses the JSON response, and makes pipeline pass/fail decisions based on the ALLOW or BLOCK verdict.

\vspace{0.3cm}
\noindent\textbf{Git Repository — Gitea (Secondary Actor)} \\
Gitea serves as the local Git server that stores infrastructure code, triggers webhooks on push events, and notifies Jenkins of new commits. It provides version control and change tracking for the developer workflow.

% ============================================================
\subsection{Functional Requirements}
% ============================================================

The functional requirements describe the expected actions from each actor within the system. These functionalities address the needs for security scanning, monitoring, and pipeline integration.

\vspace{0.3cm}
\noindent\textbf{For the Developer:}
\begin{itemize}
    \item Push infrastructure code (Terraform, Dockerfile, Kubernetes YAML) to a Git repository.
    \item Receive automated security scan results through the Jenkins pipeline console.
    \item View detailed violation messages with line numbers and fix recommendations.
    \item Access HTML security reports generated per scan.
\end{itemize}

\noindent\textbf{For the DevOps Engineer:}
\begin{itemize}
    \item Configure the block threshold (CRITICAL, HIGH, MEDIUM, LOW) to control pipeline behavior.
    \item Monitor real-time metrics including total scans, block rate, and top violations.
    \item Review scan history with timestamps, filenames, and violation counts.
    \item Access the real-time dashboard with statistics and visualizations.
    \item Generate and view HTML security reports for any scan.
\end{itemize}

\noindent\textbf{For the Security Engineer:}
\begin{itemize}
    \item Enable or disable specific security policies via configuration.
    \item Adjust severity levels for individual policies.
    \item Define required tags and general scanning parameters.
    \item Monitor violation trends and top-triggered rules.
\end{itemize}

\noindent\textbf{For Jenkins / CI System:}
\begin{itemize}
    \item Automatically detect code changes via Git polling or webhooks.
    \item Submit infrastructure code for analysis via the REST API.
    \item Parse ALLOW/BLOCK decisions and control pipeline flow accordingly.
    \item Log scan results and violation details in the build console.
\end{itemize}

% ============================================================
\subsection{Non-Functional Requirements}
% ============================================================

The non-functional requirements describe the qualitative criteria necessary to guarantee a smooth experience and optimal system operation.

\begin{itemize}
    \item \textbf{Performance}: The service must respond in less than 50 milliseconds on average to avoid creating bottlenecks in CI/CD pipelines. The target throughput is 100+ requests per second.
    \item \textbf{Reliability}: The service must produce consistent results for identical inputs. Stateless design ensures no data loss on restart.
    \item \textbf{Security}: All API endpoints must validate input data. CORS middleware must be configured. No sensitive data should be logged.
    \item \textbf{Scalability}: The architecture must support the addition of new security policies without modifying existing code. New IaC formats must be addable through the parser/normalizer pattern.
    \item \textbf{Maintainability}: The code must be structured and documented, following the 4-layer architecture. Each layer must be independently testable.
    \item \textbf{Portability}: The service must run completely offline without cloud dependencies. Docker containerization ensures cross-platform compatibility.
    \item \textbf{Usability}: Violation messages must be educational, explaining the risk and providing actionable remediation steps.
\end{itemize}

% ============================================================
\section{Product Backlog}
% ============================================================

As Product Owner, we translated the documented requirements into a structured Product Backlog, where each functionality is written as User Stories to ensure business value. A rigorous prioritization technique was applied, hierarchizing elements according to their strategic importance.

\begin{longtable}{|p{3.5cm}|p{7cm}|c|}
\caption{Product Backlog — User Stories with Priorities} \label{tab:backlog} \\
\hline
\textbf{Feature Group} & \textbf{User Story} & \textbf{Priority} \\
\hline
\endfirsthead
\hline
\textbf{Feature Group} & \textbf{User Story} & \textbf{Priority} \\
\hline
\endhead
\hline
\endfoot

Infrastructure Parsing
& As a DevOps engineer, I can submit Terraform HCL code for automated security scanning.
& High \\
\hline

& As a DevOps engineer, I can submit Dockerfile content for security analysis.
& High \\
\hline

& As a DevOps engineer, I can submit Kubernetes YAML manifests for policy evaluation.
& High \\
\hline

Format Normalization
& As a developer, I can expect consistent violation reporting regardless of the infrastructure format used.
& High \\
\hline

Security Policy Engine
& As a security engineer, I can enforce 100 security policies covering AWS, Docker, Kubernetes, and Terraform.
& High \\
\hline

& As a security engineer, I can configure the severity level (CRITICAL, HIGH, MEDIUM, LOW) at which deployments are blocked.
& High \\
\hline

& As a developer, I receive detailed violation messages explaining the risk and how to fix it.
& High \\
\hline

CI/CD Integration
& As a DevOps engineer, I can integrate the service into a Jenkins pipeline to automatically scan code on every push.
& High \\
\hline

& As a DevOps engineer, I can configure the service to scan only changed lines using Git diff.
& Medium \\
\hline

Dashboard \& Reporting
& As a DevOps engineer, I can view a real-time dashboard with scan statistics, top violations, and block rates.
& Medium \\
\hline

& As a DevOps engineer, I can generate HTML security reports for each scan.
& Medium \\
\hline

Scan History
& As a DevOps engineer, I can view the history of all past scans with timestamps and results.
& Medium \\
\hline

& As a DevOps engineer, I can access metrics including total requests, blocks, allows, and top violations.
& Medium \\
\hline

Health \& Monitoring
& As an administrator, I can check the service health status via the health endpoint.
& Low \\
\hline

Configuration
& As a security engineer, I can enable or disable individual policies through a YAML configuration file.
& Medium \\
\hline

Docker Deployment
& As a DevOps engineer, I can deploy the service using Docker and docker-compose with a single command.
& Medium \\
\hline

\end{longtable}

% ============================================================
\section{Diagrams}
% ============================================================

\subsection{Global Use Case Diagram}

The use case diagram is an essential tool for visualizing the interactions between different actors and the system. It provides an overview of the actions that users (Developer, DevOps Engineer, Security Engineer) and systems (Jenkins, Gitea) can perform. This diagram highlights the main functionalities of the system and illustrates how each actor uses them. It facilitates not only understanding of the exchanges between the system and its users, but also helps detect potential insufficiencies or inconsistencies in the overall operation.

% TODO: CREATE USE CASE DIAGRAM IN DRAW.IO
% The diagram should show:
% - Developer: Push code, View violations, View reports
% - DevOps Engineer: Configure threshold, View dashboard, View metrics, View history, Deploy service
% - Security Engineer: Configure policies, Adjust severity, Monitor trends
% - Jenkins: Trigger scan (POST /analyze), Parse response (ALLOW/BLOCK), Control pipeline
% - Gitea: Store code, Notify Jenkins
% - System: Analyze code, Generate report, Save to history, Return decision
%
% \begin{figure}[H]
%     \centering
%     \includegraphics[width=\textwidth]{figures/use-case-diagram.png}
%     \caption{Global use case diagram}
%     \label{fig:use-case}
% \end{figure}

\textit{[Figure: Insert your Use Case Diagram here — create it in draw.io following the actors and use cases described above]}

\subsection{Global Class Diagram}

The class diagram describes the structure of our system by showing the classes, their attributes and methods, as well as the relationships between them. It serves as a blueprint for code implementation, helping to understand how different parts of the system interact. The class diagram helps ensure that our design is solid, logical, and well-organized, facilitating the development phase.

% TODO: CREATE CLASS DIAGRAM IN DRAW.IO
% The diagram should show these classes and their relationships:
%
% AnalyzeRequest (code_type, content, filename, block_threshold, diff_only)
% AnalyzeResponse (decision, violations[], explanation, scan_time_ms, filename)
% Violation (rule, severity, line, message, recommendation)
%
% PolicyEngine (normalizer, policies[]) → analyze()
% Normalizer → normalize() returns NormalizedCode
% NormalizedCode (code_type, resources[], variables, metadata)
% Resource (resource_type, resource_name, attributes, line_number)
%
% BasePolicy (rule_id, severity, description) → check()
%   └── PublicS3BucketPolicy, HardcodedSecretPolicy, ... (100 policies)
%
% TerraformParser → parse(), extract_resources()
% DockerfileParser → parse()
% KubernetesParser → parse(), extract_containers()
%
% PolicyDecision (allow, violations[], summary, scan_duration_ms)
% PolicyViolation (rule_id, severity, message, line_number, recommendation)
%
% \begin{figure}[H]
%     \centering
%     \includegraphics[width=\textwidth]{figures/class-diagram.png}
%     \caption{Global class diagram}
%     \label{fig:class-diagram}
% \end{figure}

\textit{[Figure: Insert your Class Diagram here — create it in draw.io following the class structure described above]}

% ============================================================
\section{Application Architecture}
% ============================================================

In our project, we adopted a \textbf{4-layer architecture} designed to ensure clear separation of concerns, maintainability, and extensibility. Each layer has a single, well-defined responsibility, and communication flows strictly from the top layer (API) down to the bottom layer (Policy Engine).

\vspace{0.5cm}
\noindent\textbf{Layer 1 — API Layer (FastAPI)}

The API layer serves as the entry point for all external interactions. Built with FastAPI, it exposes RESTful endpoints for code analysis, metrics retrieval, health checking, scan history, and report generation. It handles input validation through Pydantic schemas, manages CORS middleware for cross-origin requests, and serializes responses to JSON format. The endpoints are:

\begin{itemize}
    \item \texttt{POST /analyze} — Submit infrastructure code for security scanning
    \item \texttt{GET /metrics} — Retrieve aggregated scan statistics
    \item \texttt{GET /health} — Service health check
    \item \texttt{GET /history} — Retrieve scan history from SQLite
    \item \texttt{GET /report/\{filename\}} — Generate HTML security report
    \item \texttt{GET /} — Serve the real-time dashboard UI
\end{itemize}

\vspace{0.5cm}
\noindent\textbf{Layer 2 — Parser Layer}

The parser layer is responsible for converting raw infrastructure code into structured data. Three specialized parsers handle the three supported formats:

\begin{itemize}
    \item \textbf{TerraformParser}: Uses the \texttt{python-hcl2} library to parse HashiCorp Configuration Language into Python dictionaries, then extracts individual resources with their types, names, and attributes.
    \item \textbf{DockerfileParser}: A custom regex-based parser that processes Dockerfile instructions line by line, extracting each directive (FROM, RUN, COPY, USER, etc.) with its arguments and line number.
    \item \textbf{KubernetesParser}: Uses the \texttt{PyYAML} library to parse YAML manifests, then extracts pods, containers, and security contexts from Deployment and Pod specifications.
\end{itemize}

\vspace{0.5cm}
\noindent\textbf{Layer 3 — Normalizer Layer}

The normalizer layer is the \textbf{key architectural innovation} of this project. It converts the heterogeneous output of the three parsers into a unified \texttt{Resource} model:

\begin{lstlisting}[language=Python, caption={Unified Resource model}, label={lst:resource}]
Resource(
    resource_type = "aws_s3_bucket",
    resource_name = "my-bucket",
    attributes    = {"acl": "public-read"},
    line_number   = 5
)
\end{lstlisting}

This normalization is critical because it decouples policies from formats:
\begin{itemize}
    \item \textbf{Without normalizer}: 100 policies $\times$ 3 formats = 300 implementations
    \item \textbf{With normalizer}: 100 policies $\times$ 1 format = 100 implementations
    \item \textbf{Result}: 67\% reduction in code complexity
\end{itemize}

\vspace{0.5cm}
\noindent\textbf{Layer 4 — Policy Engine}

The policy engine is the core decision-making component. It receives normalized code from the normalizer, runs all 100 security policies against the resources, aggregates violations, applies the severity threshold filter, and produces a final ALLOW or BLOCK decision with a detailed summary.

The 100 policies are organized into 14 categories:

\begin{table}[H]
\centering
\caption{Security policy categories and counts}
\label{tab:policy-categories}
\small
\begin{tabular}{|l|c|l|}
\hline
\textbf{Category} & \textbf{Count} & \textbf{Examples} \\
\hline
Core Policies & 12 & Hardcoded secrets, public S3, public DB \\
\hline
AWS S3 & 5 & Versioning, MFA delete, lifecycle \\
\hline
AWS EC2/Compute & 7 & IMDSv2, public IP, EBS encryption \\
\hline
AWS RDS/Database & 7 & Multi-AZ, backup retention, deletion protection \\
\hline
AWS IAM & 7 & Admin policy, wildcard, MFA, key rotation \\
\hline
Networking/VPC & 6 & Flow logs, SSH open, RDP open \\
\hline
Docker & 8 & Latest tag, healthcheck, secrets in ENV \\
\hline
Kubernetes & 12 & Readonly root FS, privilege escalation \\
\hline
CloudTrail/Monitoring & 8 & CloudTrail, GuardDuty, Security Hub \\
\hline
ELB/WAF/Route53 & 7 & HTTP listener, WAF, TLS policy \\
\hline
Terraform General & 5 & Remote backend, provider version \\
\hline
EKS/Lambda/ECS & 7 & EKS public endpoint, Lambda secrets \\
\hline
Messaging/Cache & 5 & SNS encryption, DynamoDB encryption \\
\hline
ECR/Secrets/CloudFront & 4 & ECR scanning, Secrets Manager rotation \\
\hline
\textbf{Total} & \textbf{100} & \\
\hline
\end{tabular}
\end{table}

% TODO: CREATE 4-LAYER ARCHITECTURE DIAGRAM IN DRAW.IO
% Show the 4 layers as stacked boxes with arrows:
% [API Layer (FastAPI)] → [Parser Layer (TF/Docker/K8s)] → [Normalizer (→ Resource model)] → [Policy Engine (100 policies → ALLOW/BLOCK)]
% Also show: SQLite for history, Dashboard HTML, and external actors (Jenkins, Gitea, Developer)
%
% \begin{figure}[H]
%     \centering
%     \includegraphics[width=\textwidth]{figures/architecture.png}
%     \caption{Proposed 4-layer architecture for the platform}
%     \label{fig:architecture}
% \end{figure}

\textit{[Figure: Insert your 4-Layer Architecture Diagram here — create it in draw.io]}

% ============================================================
\section{Tools and Development Environment}
% ============================================================

\subsection{Hardware Environment}

The development and testing of this project were carried out using the following hardware:

\begin{itemize}
    \item \textbf{Host Machine (Lenovo IdeaPad Gaming 3):}
    \begin{itemize}
        \item Processor: AMD Ryzen 7
        \item Memory: 16 GB RAM
        \item Storage: 500 GB SSD
    \end{itemize}
    \item \textbf{Development Virtual Machine (Ubuntu Linux):}
    \begin{itemize}
        \item Allocated Memory: 6 GB RAM
        \item Allocated CPUs: 4 cores
        \item Allocated Storage: 50 GB
        \item Operating System: Ubuntu Linux
    \end{itemize}
\end{itemize}

% ============================================================
\subsection{Technical Environment}
% ============================================================

The following technologies were used throughout the development of the project:

\vspace{0.3cm}

% TODO: For each tool below, you can optionally add its logo image.
% Download logos and place them in figures/ folder.

\noindent\textbf{Python 3.10} — The primary programming language used for all backend development. Python's rich ecosystem of libraries, readability, and rapid development capabilities made it ideal for this project.

\vspace{0.2cm}
\noindent\textbf{FastAPI 0.104.1} — A modern, high-performance web framework for building APIs with Python. FastAPI provides automatic data validation via Pydantic, interactive API documentation (Swagger UI), and native async support.

\vspace{0.2cm}
\noindent\textbf{Pydantic 2.5.0} — A data validation library that ensures type safety and automatic serialization/deserialization of API request and response models.

\vspace{0.2cm}
\noindent\textbf{python-hcl2 4.3.2} — A Python library for parsing HashiCorp Configuration Language (HCL), used to process Terraform files into Python dictionaries.

\vspace{0.2cm}
\noindent\textbf{PyYAML 6.0.1} — A YAML parser and emitter for Python, used to process Kubernetes manifest files.

\vspace{0.2cm}
\noindent\textbf{SQLite} — A lightweight, serverless relational database used to persist scan history across service restarts.

\vspace{0.2cm}
\noindent\textbf{Docker \& Docker Compose} \cite{docker} — Containerization platform used to package the service into a portable, reproducible container. Docker Compose orchestrates multi-container deployments.

\vspace{0.2cm}
\noindent\textbf{Jenkins 2.x} \cite{jenkins} — Open-source automation server used for implementing CI/CD pipelines that automatically trigger security scans on code push events.

\vspace{0.2cm}
\noindent\textbf{Gitea} — A lightweight, self-hosted Git service used as a local Git server for the offline demonstration, simulating a company-grade source control environment.

\vspace{0.2cm}
\noindent\textbf{pytest 7.4.3 \& httpx 0.25.1} — Testing frameworks used for unit testing and HTTP-based integration testing of the API endpoints.

\vspace{0.2cm}
\noindent\textbf{Visual Studio Code} \cite{vscode} — The primary code editor used for development, with extensions for Python, Docker, and YAML support.

\vspace{0.2cm}
\noindent\textbf{GitHub} \cite{github} — Used for hosting the project's source code repository (\texttt{ridha-gn/devsecops-policy-service}).

% ============================================================
\section{Conclusion}
% ============================================================

In conclusion, this chapter established the solid foundations of our project by defining the actors, functional and non-functional requirements, and by establishing a detailed Product Backlog. The UML diagrams provided a clear understanding of the application's architecture. The choice of appropriate technologies was determinant in guaranteeing a performant and reliable application. The tools and working environment selected reinforced the efficiency of our development process. This positions us well for the implementation stages that are discussed in the General Conclusion, which presents the results achieved across our four development sprints.
